% In this file you should put the actual content of the blueprint.
% It will be used both by the web and the print version.
% It should *not* include the \begin{document}
%
% If you want to split the blueprint content into several files then
% the current file can be a simple sequence of \input. Otherwise It
% can start with a \section or \chapter for instance.

\chapter{Introduction}

This project formalizes the definition of \Ainf{}-categories and constructs an
\Ainf{}-endofunctor on a variant of the KLRW category, which acts as a braid group invariant.

\chapter{\Ainf{}-algebras}

\section{Overview}

We define an \Ainf{}-algebra, which can be thought of as an \Ainf{}-category with a single element.
This will later be generalized with the definition of an \Ainf{}-category.

\section{Gradings}

\begin{definition}\label{GradingIndex}
  A \emph{grading index} is an abelian group $(\beta,0,+)$ equipped with additive group
  homomorphisms $\phi : \mathbb{Z} \to \beta$ and $\sigma : \beta \to \mathbb{Z}/2\mathbb{Z}$ such
  that $\sigma \circ \phi$ is the canonical projection $\pi : \mathbb{Z} \to \mathbb{Z}/2\mathbb{Z}$.

\end{definition}

\begin{definition}[Mathlib]\label{GradedObject}
  Let $\mathcal C$ be a category and $\beta$ a type (not necessarily a grading index).
  Then a \emph{$\beta$-graded object in $\mathcal C$} is a function $X : \beta \to \mathcal C$.
  The collection of $\beta$-graded objects in $\mathcal C$ forms a category $\mathrm{Gr}_\beta(\mathcal C)$.

\end{definition}

\begin{definition}\label{GradedModule}

  Given a grading index $\beta$ and a commutative ring $R$, a \emph{graded $R$-module} is a
  $\beta$-graded object in the category $R\text{-}\mathrm{Mod}$.
\end{definition}
